\DeclareEditionOverlay{CM-C04-U001}{%
\begin{EditionUnitCard}{La palabra estrictamente importa}
Compara las desigualdades de entrada y salida en cada definición.
\begin{EditionCheckpointBox}{Comprueba}
Da un ejemplo creciente que no sea estrictamente creciente.
\end{EditionCheckpointBox}
\end{EditionUnitCard}}

\DeclareEditionOverlay{CM-C04-U002}{%
\begin{EditionUnitCard}{Recupera las herramientas}
Vuelve a los enunciados del capítulo anterior y localiza sus hipótesis.

\end{EditionUnitCard}}

\DeclareEditionOverlay{CM-C04-U003}{%
\begin{EditionUnitCard}{Del signo a la variación}
Un intervalo de signo constante permite una conclusión en todo el intervalo.
\begin{EditionCheckpointBox}{Comprueba}
Indica el intervalo sobre el que conoces el signo antes de concluir monotonía.
\end{EditionCheckpointBox}
\end{EditionUnitCard}}

\DeclareEditionOverlay{CM-C04-U004}{%
\begin{EditionUnitCard}{Conclusiones sobre intervalos}
La conexión del dominio importa al pasar de derivadas a funciones.
\begin{EditionWarningBox}{Atención}
No saltes entre componentes separadas del dominio.
\end{EditionWarningBox}
\end{EditionUnitCard}}

\DeclareEditionOverlay{CM-C04-U005}{%
\begin{EditionUnitCard}{Control local, efecto global}
Una cota de la derivada limita cuánto puede cambiar la función.
\begin{EditionCheckpointBox}{Comprueba}
Identifica la cota y la distancia entre los puntos.
\end{EditionCheckpointBox}
\end{EditionUnitCard}}

\DeclareEditionOverlay{CM-C04-U006}{%
\begin{EditionUnitCard}{Candidato no significa extremo}
Construye una lista completa y después clasifica.
\begin{EditionWarningBox}{Atención}
Incluye puntos donde no existe la derivada y revisa los extremos del intervalo.
\end{EditionWarningBox}
\end{EditionUnitCard}}

\DeclareEditionOverlay{CM-C04-U007}{%
\begin{EditionUnitCard}{Lee el cambio de signo}
El comportamiento a ambos lados decide el tipo de extremo.
\RegisteredBookFigure{FIG-C04-001}{fig:c04:derivative-sign}{Del signo de la derivada a la variación.}{Un cuadro de signos muestra crecimiento, decrecimiento y crecimiento, con un máximo local y un mínimo local en los cambios de signo.}{\DerivativeSignVisual}
\begin{EditionCheckpointBox}{Comprueba}
Completa un cuadro de signos antes de nombrar máximo o mínimo.
\end{EditionCheckpointBox}
\end{EditionUnitCard}}

\DeclareEditionOverlay{CM-C04-U008}{%
\begin{EditionUnitCard}{Cero puede ser inconcluso}
El criterio rápido funciona solo bajo sus hipótesis.
\begin{EditionWarningBox}{Atención}
Si f''(a)=0, vuelve a un criterio que examine el comportamiento alrededor.
\end{EditionWarningBox}
\end{EditionUnitCard}}

\DeclareEditionOverlay{CM-C04-U009}{%
\begin{EditionUnitCard}{Existencia y comparación}
Weierstrass garantiza; la derivada localiza candidatos; la comparación decide.
\begin{EditionCheckpointBox}{Comprueba}
Haz una tabla con todos los candidatos y los extremos del intervalo.
\end{EditionCheckpointBox}
\end{EditionUnitCard}}

\DeclareEditionOverlay{CM-C04-U010}{%
\begin{EditionUnitCard}{Modelar antes de derivar}
Nombra variables y restricciones antes de construir la función objetivo.
\begin{EditionCheckpointBox}{Comprueba}
Escribe qué valores son admisibles y comprueba la respuesta en el contexto.
\end{EditionCheckpointBox}
\end{EditionUnitCard}}

\DeclareEditionOverlay{CM-C04-U011}{%
\begin{EditionUnitCard}{Empieza por la geometría}
Usa la posición respecto de las cuerdas para fijar la convención.
\begin{EditionWarningBox}{Atención}
No cambies de convexo a cóncavo por una convención memorizada de otra fuente.
\end{EditionWarningBox}
\end{EditionUnitCard}}

\DeclareEditionOverlay{CM-C04-U012}{%
\begin{EditionUnitCard}{Tres lenguajes de convexidad}
Relaciona cuerdas, crecimiento de la derivada y signo de la segunda derivada.
\begin{EditionCheckpointBox}{Comprueba}
Anota las hipótesis necesarias para pasar de un lenguaje a otro.
\end{EditionCheckpointBox}
\end{EditionUnitCard}}

\DeclareEditionOverlay{CM-C04-U013}{%
\begin{EditionUnitCard}{Globalidad no garantiza existencia}
La convexidad ayuda a clasificar y asegurar unicidad si existe el candidato adecuado.
\begin{EditionWarningBox}{Atención}
Separa las preguntas: existe, es global, es único.
\end{EditionWarningBox}
\end{EditionUnitCard}}

\DeclareEditionOverlay{CM-C04-U014}{%
\begin{EditionUnitCard}{La inflexión exige cambio}
Examina la convexidad a ambos lados del punto.
\begin{EditionWarningBox}{Atención}
f''(a)=0 por sí solo no demuestra una inflexión.
\end{EditionWarningBox}
\end{EditionUnitCard}}

\DeclareEditionOverlay{CM-C04-U015}{%
\begin{EditionUnitCard}{Ampliación: una propiedad de las derivadas}
Lee el resultado como una conexión, no como una nueva técnica obligatoria.

\end{EditionUnitCard}}

